% Do not change document class, margins, fonts, etc.
\documentclass[a4paper,oneside,bibliography=totoc]{scrartcl}

% some useful packages (add more as needed)
\usepackage[utf8]{inputenc}
\usepackage[T1]{fontenc}
\usepackage{graphicx}
\usepackage{latexsym}
\usepackage{amsmath}
\usepackage{amssymb}
\usepackage{tabularx}
\usepackage{booktabs}
\usepackage{enumitem}
\usepackage{algorithm} % you can modify the algorithm style to your liking
\usepackage{algorithmic}
\usepackage{csquotes}
\renewcommand{\algorithmiccomment}[1]{\hfill\textit{// #1}}
\usepackage[usenames,dvipsnames]{xcolor}
\newcommand{\todo}[1]{\textcolor{red}{\textbf{[TODO: #1]}}}
\usepackage[colorlinks,citecolor=Green]{hyperref} % you may change/remove the colors
\usepackage{tcolorbox}
\tcbuselibrary{skins,breakable}
\usepackage{longtable}
\usepackage{microtype}

% chicago citation style
\usepackage{natbib}
\bibliographystyle{chicagoa}
\setcitestyle{authoryear,round,semicolon,aysep={},yysep={,}} \let\cite\citep

% example enviroments (add more as needed)
\newtheorem{definition}{Definition} \newtheorem{proposition}{Proposition}

\newtcolorbox{instructionbox}{
  colback=white,
  colframe=black,
  boxrule=0.8pt,
  arc=0pt,
  breakable,
  enhanced,
  fontupper=\small
}

\begin{document}

\subject{Computational Analysis of Communication}
\title{Measuring the Crime Frame in German News Media}
\author{Robin-Kiara Braun \and Marmee Pandya}
\date{\today}
\maketitle

\begin{abstract}
Racism in the news rarely announces itself all at once. It accumulates, paragraph by paragraph, in small choices about which crimes get a nationality attached to them and which don't. Whether a German newspaper names a group at all when reporting an offense, and whether that mention reads as an isolated fact or as evidence of a group's inherent danger, is a framing choice with documented effects on public attitudes toward immigration and minority groups \citep{devreese2011,guo2023,lecheler2015}. Measuring that choice across an entire national news cycle rather than in a hand-coded sample of a few hundred articles has been difficult. We build and validate a large language model (LLM) annotation pipeline that does exactly this, across 659{,}895 German news paragraphs from 60 outlets spanning 78 ethnic, national, and religious groups in 2022, grounded in \citet{freudenthaler2025}'s framework for measuring fragmented racism in large-scale corpora. We operationalize the security-crime frame as a binary judgment, \texttt{CRIME\_FRAME\_NEG} versus \texttt{NO\_CRIME\_FRAME}, having started from a three-class scheme that also distinguished positively-valenced crime framing and collapsed it once positive framing turned out to occur in roughly even less than one percent of paragraphs, too rare for either human coders or LLMs to evaluate reliably. Getting the remaining binary judgment right required resolving genuinely hard interpretive questions that recur throughout German crime reporting: does suspicion count before conviction? Does a crime committed abroad implicate a group at home? Does legitimizing violence count the same as committing it? We show that resolving such questions explicitly, and revising them against observed disagreement, is what separates a defensible measurement instrument from an arbitrary one: human intercoder reliability rose from Krippendorff's $\alpha = 0.56$ to $0.79$ across two independently drawn 250-paragraph batches as the codebook matured, and LLM validity against human ground truth reached $\kappa = 0.70$ ($F_1 = 0.85$) on a held-out, independently drawn validation batch under a revised few-shot prompting strategy evidence that these gains reflect a genuinely better-specified instrument rather than overfitting to a single test set. \todo{need to update it based on the 600 entries result}
\end{abstract}

\section{Introduction}
\label{sec:intro}

News media do not simply report on ethnic, national, and religious groups; they select which aspects of a group's presence to foreground, and threat is one of the most consequential of these selections. When a newspaper links a group to crime, danger, or violence, it activates what \citet{entman1993} calls a frame: a way of ``selecting some aspects of a perceived reality and making them more salient \dots\ in such a way as to promote a particular problem definition.'' Framing effects research has shown that exposure to such threat frames measurably shifts audience emotions and policy attitudes toward immigrants and minority groups \citep{devreese2011,lecheler2015,guo2023}, and that these threat and benefit portrayals are not evenly distributed across corpora \citep{eberl2018}.

This project operationalizes and measures one such frame at scale: the \emph{security-crime frame}, understood as the degree to which a news paragraph portrays a group as a source of criminal or security threat (negative valence) versus a contributor to safety, protection, or crime prevention (positive valence), versus containing no such frame at all. We build this measurement on top of \citet{freudenthaler2025}'s recently proposed framework for measuring fragmented traces of racism in large-scale news corpora, which places threat and benefit framing alongside five other sub-dimensions of racist and anti-racist discourse (social position, entitlement, essentialism, homogenization, and othering) as one of the elements that link ideology to structural inequality and audience attitudes (see Section~\ref{sec:theory}).

% 2022 is a substantively motivated choice of year for this question, not an arbitrary one: German news coverage of migration and refugee movements changed markedly following the outbreak of the war in Ukraine in February 2022, at the same time as long-running domestic debates about asylum policy, integration, and crime reporting continued. A corpus spanning this year lets us observe security-crime framing applied to newly prominent groups (Ukrainian refugees), long-discussed ones (asylum seekers, Muslims, various other nationalities), and groups comparatively rarely framed as security-relevant at all (Western European and North American nationalities), all within a single, coherent news cycle, a precondition for the cross-group comparison our research question asks for.
\todo{do we want to add something about the choice of 2022 as the year?}

Two practical problems make this measurement difficult at scale. First, the underlying corpus is large: our dataset comprises 659{,}895 German-language news paragraphs from 60 outlets published in 2022, mentioning 78 distinct ethnic, national, or religious groups, far too many for exhaustive human coding. Second, framing is a subtle, context-dependent judgment that simple keyword counting cannot capture reliably: the same lexical item (e.g.\ ``crime,'' ``police,'' ``war'') can appear in paragraphs that clearly do and clearly do not frame a group as a security threat, depending on whether the group is cast as perpetrator, victim, bystander, or unrelated third party. We therefore rely on large language models (LLMs) prompted with detailed natural-language coding instructions, following the human-LLM hybrid annotation. Therefore, we (1) draft instructions, (2) test and improve their coverage, (3) validate against a curated sample of hard and soft cases, and (4) validate against a representative, independently human-coded sample.

% \emph{How is the security-crime frame differentially applied to ethnic, national, and religious groups in German news coverage, and does this vary systematically by outlet type and political orientation?}
Our research question is: \todo{research question goes here}  Answering this requires, as a precondition, an annotation pipeline whose reliability and validity can be reported with the same rigor as the substantive findings it produces. \todo{if we can do that it would be so cool but im not sure if we have enough git history :( } Accordingly, this report makes two contributions. First, we document, transparently and in more detail than is typical for the ``we improved the prompt in an iterative process'' sentence that most annotation studies settle for, the full evolution of an LLM-based annotation pipeline: what we changed between iterations, why, and what did and did not work, including a negative result on confidence-based model routing and a positive result on multi-model ensembling. Second, we report on the framing patterns the validated pipeline reveals across groups and outlet types. % TODO: once n=600 results and group/outlet crosstab are available, extend this paragraph with the headline substantive result.


Concretely, this report contributes:
\begin{itemize}[noitemsep]
    \item a fully specified, iteratively refined LLM annotation scheme for the security-crime frame, including the decision rules that separate it from adjacent concerns such as war reporting, fiction, and generic political protest (Section~\ref{sec:method}, Appendix~\ref{app:instructions});
    \item a transparent account of four annotation-development stages, including two side-experiments (multi-model ensembling and confidence-based routing) with one positive and one negative result (Sections~\ref{sec:method-ensemble}--\ref{sec:method-confidence});
    \item human intercoder reliability and LLM validity estimates at each stage, including a diagnosed discrepancy between rising human reliability and non-improving model validity across batches (Section~\ref{sec:results});
    \item and, building on the validated pipeline, a substantive account of how the security-crime frame is distributed across ethnic, national, and religious groups and across outlet types in our corpus (Section~\ref{sec:results-substantive}).
\end{itemize}

The remainder of this report is organized as follows. Section~\ref{sec:theory} situates the security-crime frame within \citet{freudenthaler2025}'s broader framework and the wider framing-effects literature. Section~\ref{sec:data} describes the corpus, its outlets, and the sampling and preprocessing steps. Section~\ref{sec:method} documents the four-stage annotation pipeline, the models used, and the two side-experiments we ran (LLM ensembling and confidence-based routing). Section~\ref{sec:results} reports intercoder reliability, LLM validity, and the substantive framing patterns we find. Section~\ref{sec:discussion} discusses what these results mean, Section~\ref{sec:limitations} discusses limitations and future work, and Section~\ref{sec:conclusion} concludes.

\section{Theoretical Background}
\label{sec:theory}

\subsection{Racism as a Fragmented Discourse Pattern}

Racism research distinguishes between the structural persistence of racial inequality and the ideological, attitudinal, and behavioral processes that maintain it \citep{fredrickson2002,omiwinant2014,mills2017,bonillasilva2001}. Communication research contributes to this picture because news production is one of the arenas in which racial ideology is reproduced or contested \citep{vandijk1991,feagin2020}. A recent systematic review of computational approaches to measuring racism in text finds that most studies still operationalize a single, narrow indicator rather than a multidimensional construct \citep{kathirgamalingam2025}, which is precisely the gap \citeauthor{freudenthaler2025}'s framework is designed to close. \citet{freudenthaler2025} argue that this contribution has so far been fragmented: individual studies typically capture one sub-dimension of racist discourse: the social roles groups are cast in \citep{shorvanderijt2024}, the stereotype content ascribed to them \citep{cuddy2008,fiske2018}, the degree to which they are essentialized or homogenized \citep{rothvanstee2023}, or the threat and benefit frames applied to them \citep{eberl2018,guo2023}, without integrating these dimensions into a single measurement framework. \citeauthor{freudenthaler2025}'s framework proposes six such dimensions (social position, entitlement, essentialism, homogenization, othering, and threat/benefit framing) and argues that measuring their combination, aggregated over a large corpus, lets researchers place any given group's portrayal on a spectrum from overtly racist to stereotyping, moderating, and egalitarian (or anti-racist) discourse, rather than judging individual texts as racist or not.

This project isolates and operationalizes the sixth dimension, threat and benefit framing, and narrows it further to a single frame family: the \emph{security-crime frame}. We made this scoping decision deliberately, in line with \citeauthor{freudenthaler2025}'s own observation that a full multidimensional measurement is rarely feasible within the scope of a single study, and because the security-crime frame has an established place in the framing-effects literature as one of the more behaviorally consequential frames applied to migrant and minority groups (Section~\ref{sec:framing-lit}).

\subsection{Framing Theory and the Security-Crime Frame}
\label{sec:framing-lit}

\citet{entman1993} defines framing as the selective emphasis of some aspects of an issue over others so as to promote a particular interpretation. Applied to immigration and minority-group coverage, the security-crime frame casts a group's presence primarily in terms of danger, criminality, or the erosion of public safety, as opposed to alternative frames that might emphasize a group's economic contribution, cultural difference, or humanitarian circumstances \citep{eberl2018}. \citet{devreese2011} show, in the context of news coverage of Turkish EU membership, that even a relatively small share of coverage employing a security threat frame (five percent, in their sample) was associated with measurable shifts in public support, and that negative frames were substantially more prevalent than positive ones. \citet{guo2023} similarly find that cumulative exposure to crime frames in immigration coverage shapes audience emotions and opinions over and above any single article's effect, and that such crime frames tend to produce stronger and more consistent effects across audiences than competing frames such as the culture frame. \citet{lecheler2015} show that these effects are mediated by discrete emotions, particularly anxiety, that framing evokes in the audience. This is why the valence of the security-crime frame, threat versus protection, and not merely its presence, is the theoretically important quantity to measure.

We therefore define the security-crime frame as present whenever a group is explicitly linked to crime, security, danger, violence, terrorism, suspicion, investigation, or prevention, and we further distinguish two valences: a \emph{negative} valence, in which the group is cast as a source of threat (as perpetrator, suspect, or legitimizer of violence), and a \emph{positive} valence, in which the group is cast as contributing to safety (through crime prevention, integration work, or explicit differentiation from a generalized suspicion). Paragraphs in which crime or security is discussed without an explicit link to the group, including paragraphs in which the group instead appears as a victim, paragraphs describing war or geopolitical conflict, and fictional or satirical depictions, receive no frame. The full operational definition is given in Section~\ref{sec:method} and reproduced verbatim in Appendix~\ref{app:instructions}.

\subsection{LLMs as Content-Analysis Coders}

A growing literature evaluates whether large language models can substitute for, or usefully assist, human coders in content analysis. \citet{vanatteveldt2021} find that only human coders reliably exceed a Krippendorff's $\alpha$ of 0.8 on sentiment annotation, with dictionary methods reaching only $\alpha \approx 0.34$ and deep-learning methods $\alpha \approx 0.5$, better than chance but well short of human-level agreement. \citet{macphail2016} and \citet{oconnorjoffe2020} treat intercoder reliability as a diagnostic tool rather than a technical checkbox: persistent disagreement between coders (human or machine) signals that a concept is not yet clearly enough defined in the coding instructions, and iterating on the instructions in response to observed disagreement is itself part of a defensible measurement procedure. \citet{jumle2025} apply transformer-based models directly to generic news-frame detection, illustrating that framing categories, unlike simpler sentiment or topic labels, remain a comparatively hard target for automated coding. \citet{weberreichardt2023} formalize the choice between two annotation strategies: fine-tuning an established classifier such as BERT, which requires a large existing labeled dataset, versus prompting a generative LLM directly, which is preferable when little annotated data exists but still requires the same zero-shot, one-shot, few-shot, and chain-of-thought prompting variations we test in Section~\ref{sec:method}. Their central finding, echoed throughout our own results (Section~\ref{sec:results}), is that generative LLMs cannot be used for precise annotation \enquote{out of the box}: careful validation and tailored prompt engineering are necessary regardless of which strategy is chosen. This is the backdrop against which we designed our four-stage pipeline (Section~\ref{sec:method}): rather than treating a single round of prompting as sufficient, we treat both human and LLM annotation as processes whose reliability must be measured, diagnosed, and iteratively improved, \todo{dont know if we can do that efficiently but} and we report each stage of that process rather than only the final numbers. 

\section{Data}
\label{sec:data}

\subsection{Corpus}

Our corpus consists of 659{,}895 German-language news paragraphs (the \texttt{text\_block} unit of analysis) drawn from 60 online news outlets, all published in 2022. Each paragraph is associated with a masked group identifier (the \texttt{group} column), a source outlet (\texttt{pub}), a within-article paragraph index (\texttt{par\_index}), and a character length (\texttt{length}). The corpus mentions 78 distinct ethnic, national, and religious groups; the ten most frequent are refugees generically (\texttt{REFUG}, 117{,}132 paragraphs), Ukrainians (72{,}028), Russians (70{,}398), British (37{,}580), Americans (35{,}357), Hungarians (34{,}563), French (27{,}939), migrants generically (\texttt{MIGR}, 25{,}207), Jewish people (\texttt{JUDE}, 24{,}779), and Spanish (15{,}249). This mix of country-of-origin categories and broader social categories (refugees, migrants, Muslims, Kurds) reflects how German news coverage itself categorizes group membership, rather than a categorization scheme imposed by the research team.

The 78 groups fall into two broad categories: country-of-origin nationalities (e.g.\ \texttt{UKR}, \texttt{RUS}, \texttt{FRA}, \texttt{TUR}, \texttt{AFG}), which make up the majority of distinct group codes, and supranational, religious, or generic social categories that cut across nationality (e.g.\ \texttt{REFUG} for refugees generically, \texttt{MIGR} for migrants generically, \texttt{MUSL} for Muslims, \texttt{JUDE} for Jewish people, \texttt{KURD} for Kurds, \texttt{ASYL} for asylum seekers). This distinction matters for interpreting any group-level framing comparison (Section~\ref{sec:results-substantive}): generic categories such as \texttt{REFUG} and \texttt{MIGR} aggregate coverage across many underlying nationalities and are, by construction, more likely to co-occur with policy and security discourse than a specific nationality code is, simply because they are the categories German news coverage itself reaches for when discussing migration policy in the abstract rather than a specific country's citizens.

Outlets are classified into four categories, following a shared outlet list maintained across the seminar's research groups (Table~\ref{tab:outlets}): \emph{mainstream national} outlets (e.g.\ \textit{S\"uddeutsche Zeitung}, \textit{FAZ}, \textit{Spiegel}, \textit{Bild}), \emph{regional} outlets (e.g.\ \textit{Tagesspiegel}, \textit{Merkur}, \textit{Weser Kurier}), and two categories of \emph{alternative} media distinguished by political orientation, \emph{alternative right} (e.g.\ \textit{Epoch Times}, \textit{Compact}, \textit{PI News}) and \emph{alternative left} (e.g.\ \textit{Junge Welt}, \textit{Telepolis}, \textit{Nachdenkseiten}). Mainstream and regional outlets together account for roughly nine in ten paragraphs in the corpus, with alternative outlets, despite their much smaller share of overall coverage, providing the variation in political orientation needed to test whether framing differs by outlet ideology.

\begin{table}[htbp]
\centering
\begin{tabular}{lrrr}
\toprule
\textbf{Outlet category} & \textbf{Outlets} & \textbf{Paragraphs} & \textbf{Share} \\
\midrule
Mainstream national & 14 & 333{,}782 & 50.6\% \\
Regional            & 17 & 259{,}678 & 39.4\% \\
Alternative right    & 16 & 41{,}186  & 6.2\%  \\
Alternative left     & 13 & 25{,}249  & 3.8\%  \\
\midrule
Total                & 60 & 659{,}895 & 100.0\% \\
\bottomrule
\end{tabular}
\caption{Corpus composition by outlet category.}
\label{tab:outlets}
\end{table}

\subsection{Preprocessing and Sampling}

The corpus was originally scraped in German, and our LLM annotators are prompted in German throughout Steps 2--4 (Section~\ref{sec:method}), operating directly on the original German text. We drop paragraphs shorter than 15 words, as these were found, during initial coverage testing, to rarely contain enough context to support a framing judgment either for human or LLM coders. Group identities are masked in the text presented to annotators (as \texttt{[Gruppe~1]} or \texttt{[andere Gruppe]}, i.e.\ ``[Group~1]''/``[other group]'') to prevent both human and LLM coders from substituting a prior stereotype about a named group for a judgment grounded in the paragraph's actual content; the true group identity is reattached after annotation via the \texttt{group} column for downstream analysis. Because obvious, unambiguous security-crime framing is rare in a random sample of general news text, we oversampled candidate paragraphs during the coverage and hard-case testing stages using a zero-shot pre-classification pass, to ensure our coding instructions were tested against a reasonably balanced mix of clear and ambiguous cases rather than an overwhelming majority of unframed paragraphs (see Section~\ref{sec:method}).

\subsection{Ethical Considerations}

Measuring how groups are framed as security threats requires handling source material, including alternative-right outlets whose explicit editorial position is to frame migrant and minority groups as dangerous, a rhetorical strategy at the core of right-wing populist discourse more broadly \citep{wodak2015}, without the measurement process itself amplifying or endorsing that framing. We follow two safeguards adopted across the seminar's research groups for this reason. First, group identity is masked from annotators (both human and LLM) during coding (Section~\ref{sec:data}), so that a label is assigned based on what a paragraph explicitly says rather than on a prior stereotype the annotator might hold about a named nationality or religion; group identity is reattached only after labeling, for aggregate analysis. Second, our own reporting of results aggregates across many paragraphs and groups rather than singling out identifiable individuals mentioned in source articles, consistent with the framework's stated goal of mapping \emph{cumulative discourse patterns} rather than adjudicating whether any individual article or journalist is racist \citep{freudenthaler2025}.

\section{Method}
\label{sec:method}

We follow \citet{freudenthalerbacchuss}'s four-stage annotation-development routine for building and validating LLM-based content-analysis instructions, in which a first, non-representative round of instruction drafting, testing, and revision (Steps 1--3) precedes a final, representative human validation (Step 4). Each stage below fed directly into the next: coverage problems found in Step 2 were fixed before Step 3 began, error patterns found in Step 3 became the few-shot examples used in the instruction-revision round that follows it, and the best-performing configuration from that round, together with the model-comparison experiments described alongside it, is what we carried into Step 4.

\subsection*{Step 1: Drafting Annotation Instructions}

We drafted a three-class coding scheme: \texttt{NO\_CRIME\_FRAME}, \texttt{CRIME\_FRAME\_NEG} (the group is portrayed as a source of crime, danger, or security threat), and \texttt{CRIME\_FRAME\_POS} (a security-crime frame is present, but the group is portrayed as contributing to safety, protection, or crime prevention). This is itself an adaptation of the shared seminar codebook \citep{freudenthalercodebook}, which codes security threat and security benefit as two \emph{independent} binary flags per paragraph (each \enquote{Kommt vor: 1, Kommt nicht vor: 0}), allowing both, either, or neither to be present at once. We collapsed these two independent flags into a single, mutually exclusive three-way label to keep the annotation task tractable for a single LLM API call per paragraph.

Because small LLMs are prone to label drift, outputting a near-miss token instead of one of the three valid labels, every raw model response is parsed by matching against the label vocabulary in a fixed order (checking for \texttt{NO\_CRIME\_FRAME} before the other two, since a model that outputs the negation is otherwise at risk of also matching a laxer substring check for \texttt{CRIME\_FRAME\_NEG}). 

Several of our core exclusion rules are inherited directly from the shared seminar codebook for the security dimension of \citeauthor{freudenthaler2025}'s framework \citep{freudenthalercodebook}, which every research group in this seminar received as a common starting point and which our project narrows and adapts specifically for LLM annotation. That codebook already excludes war and geopolitical-conflict reporting, and humanitarian crises or natural disasters, from its \emph{Sicherheit} (security) threat topos, even when a group's nationality is mentioned alongside violence (e.g.\ Russian or Ukrainian military action); we adopted this exclusion into our LLM instructions unchanged. Our reasoning for why this exclusion is theoretically sound is that nationality is frequently named in war reporting for straightforwardly journalistic reasons, to identify the parties to a conflict, rather than to cast that nationality as inherently criminal or dangerous in the sense the security-crime frame is meant to capture; coding such mentions as \texttt{CRIME\_FRAME\_NEG} would conflate routine descriptive war reporting with the ideological threat-framing the measure is designed to detect. Beyond the shared codebook, we added two further exclusions specific to our narrower crime-frame operationalization: fictional, satirical, or otherwise non-factual depictions (television, film, cartoons), which are excluded regardless of content, and demonstrations or protests, which only count as framing when explicitly tied to concrete security-related consequences rather than reported as generic political expression. The full instructions, as used in the final pipeline, are reproduced in Appendix~\ref{app:instructions}.

A second inherited rule, easy to overlook but theoretically consequential, concerns undocumented migration status. The shared codebook \citep{freudenthalercodebook} explicitly states that referring to a group as \enquote{illegal immigrants} does \emph{not}, by itself, trigger the security threat topos unless an additional, separately named offense is present, while describing a group member as a smuggler (\enquote{Schleuser}) does count, since human smuggling is itself a criminal offense rather than a status description; we carried this rule into our \texttt{CRIME\_FRAME\_NEG} definition unchanged. We consider this line theoretically important regardless of its origin: had the instructions instead treated \enquote{illegal} status itself as sufficient for a crime frame, the annotation instrument would have imported, rather than measured, the very conflation of undocumented status with criminality that critical framing scholarship treats as one mechanism by which migration reporting constructs threat \citep{eberl2018}. Beyond the shared codebook, we extended \texttt{CRIME\_FRAME\_NEG} further, to \emph{ideological legitimization}: a group is coded as negatively framed when a member commits an offense, or when the group is portrayed as endorsing, celebrating, or providing ideological cover for an attack or act of violence, even absent direct involvement. This is our own addition, widening the frame from a narrow \enquote{who did it} reading to one that also captures rhetorical support for violence, which we judged closer to how threat framing plausibly operates in practice.

A third decision concerns the evidentiary threshold for \texttt{CRIME\_FRAME\_NEG}: suspicion, ongoing investigation, and unconfirmed allegations count as soon as they are explicitly linked to the group (signalled by hedges such as \enquote{mutmaßlich}/allegedly, \enquote{verdächtigt}/suspected, or \enquote{Ermittlungen gegen}/investigation against), and a judicial conviction is not required. This mirrors how crime news is conventionally reported, since the overwhelming majority of crime coverage concerns suspects rather than convicts, but it also means our measure is sensitive to \emph{allegation}, not just to adjudicated fact, which is a scope decision worth stating plainly rather than leaving implicit. Finally, for paragraphs the explicit rules above do not resolve, the instructions fall back on an audience-perception heuristic: annotators are asked to judge whether a typical reader would perceive the group as criminal, suspect, security-threatening, supportive of violence, or as a safety factor, and to default to \texttt{NO\_CRIME\_FRAME} if not. This fallback is not an afterthought: it operationalizes, at the level of a single ambiguous paragraph, the same reader-oriented definition of framing that motivates the whole project (Section~\ref{sec:framing-lit}): a frame is defined by what it makes salient to an audience \citep{entman1993}, not by an annotator's independent judgment of objective threat, so falling back on \enquote{how would a reader read this} is a theoretically motivated tie-breaker rather than an arbitrary one.


\subsection*{Step 2: Coverage and Hard-Case Testing}

We first tested whether the draft instructions covered the concept broadly, running the instructions on small samples (50 paragraphs) and visually inspecting whether clear cases were labeled sensibly and whether any recurring pattern in the data was left undefined by the instructions. We then tested for \emph{hard cases}: paragraphs on which repeated LLM annotation runs, at an elevated sampling temperature ($T=0.7$, five repeated runs per paragraph), disagreed with each other, indicating that the instructions left room for inconsistent interpretation. \todo{Here the hard cases calculation by Robin goes}

\subsection*{Step 3: Validation Against a Curated Non-Representative Sample}

Once coverage looked broad, we assembled a non-representative test set of 500 paragraphs mixing hard cases identified in Step 2, clear examples of each label, and additional keyword-searched cases, and had both authors independently hand-label a first batch of 250 of these paragraphs to serve as a human ``gold standard.'' Model annotation (\texttt{ministral-3-14b}, university API, temperature 0) against this human standard achieved 65.6\% agreement, a macro-$F_1$ of 0.521 (the unweighted average of per-class $F_1$, chosen over accuracy or micro-averaged measures specifically because it does not let a dominant class mask poor performance on rarer ones; \citealp{sokolova2009}), and a Cohen's $\kappa$ of 0.321 (Table~\ref{tab:llm-validity}, row ``Step 3''); the 86 disagreement cases were extracted as model hard cases for further instruction refinement, using consensus-coded curated samples for this stage rather than a purely random sample.

\subsection{Prompting Strategy Comparison and Instruction Revision}

This stage is where prompting strategy was systematically varied, still on the same non-representative 250-paragraph set used in Step 3 rather than a fresh sample. We compared three prompting conditions: zero-shot, few-shot (eight hand-written examples targeting error patterns observed in Step 3), and few-shot with chain-of-thought explanations.

In writing the eight few-shot examples, we followed established curation principles: an approximately balanced number of examples per label rather than mirroring the corpus's natural (heavily skewed) label distribution; at least one unambiguous, easy example per label to keep the model anchored to the core concept; and additional examples chosen specifically to exemplify a decision rule the model had gotten wrong in Step~3 (e.g.\ the victim-vs.-perpetrator distinction, and the war-reporting exclusion). We wrote these examples ourselves rather than drawing them from the corpus, to avoid contaminating the held-out evaluation set with paragraphs the model had already seen during prompting. Few-shot prompting outperformed both zero-shot and few-shot-with-CoT on this set (macro-$F_1 = 0.509$ vs.\ 0.492 and 0.502 respectively; Table~\ref{tab:llm-validity}). We then revised the instructions again and re-ran the same three conditions. Among other refinements, two unrelated rules were added in this round: foreign incidents now only count toward \texttt{CRIME\_FRAME\_NEG} when an explicit connection to Germany exists, narrowing the frame to incidents actually relevant to German society rather than any crime story naming a nationality. Separately, we clarified that deportations are not automatically \texttt{CRIME\_FRAME\_POS}, since a deportation centers the deportee as the target of an adverse state action rather than as evidence of a safety benefit. The revised few-shot condition reached accuracy $0.792$, macro-$F_1 = 0.624$, and $\kappa = 0.588$, the best result obtained on this test set (Table~\ref{tab:llm-validity}, ``Revised instructions'').

To check whether this improvement generalized beyond the specific 250 paragraphs the instructions had by then been repeatedly tuned against, we drew a second, independent batch of 250 paragraphs (``new250'') and had both authors hand-label it from scratch. Human intercoder reliability on this second batch, measured with Krippendorff's $\alpha$, the standard chance-corrected agreement statistic for nominal-category content analysis \citep{neuendorf2017}, was substantially higher than on the first ($\alpha = 0.786$ vs.\ $0.560$; Section~\ref{sec:results-reliability}), reflecting the more decisive, better-specified codebook by this point. Model validity against the new250 ground truth, however, did not improve correspondingly: macro-$F_1$ on the full 250-paragraph new250 set was $0.492$ ($\kappa = 0.498$), close to the original Step~3 result rather than the improved result from this revision round (Table~\ref{tab:llm-validity}, ``new250''). We return to this discrepancy, rising human reliability without correspondingly rising model validity, in Section~\ref{sec:discussion}, together with the note we recorded at the time flagging the possibility of instruction overfitting to the first test set.

\subsection{Comparing Different LLMs}

Because a single small model's validity plateaued around macro-$F_1 \approx 0.5$--$0.6$ despite the revised few-shot instructions above, we turned to comparing and combining different LLMs, the activity \citet{freudenthalerbacchuss}'s routine designates Step 3.4. We tested two ways of using a second, larger model, both evaluated on the same 250-paragraph test set used in Step 3 and the revision round above: first confidence-based routing (cheaper, tried first), then multi-model ensembling (more expensive, tried once routing failed).

\subsubsection{Confidence-Based Routing}
\label{sec:method-confidence}

The motivation was cost: rather than sending every paragraph through the larger, slower \texttt{Llama-3.3-70B} model, we wanted \texttt{ministral-3-14b} to flag which paragraphs it was actually unsure about and route only those onward. \texttt{ministral-3-14b} rates its own confidence in each label on a 1--5 scale, and paragraphs at or below a confidence threshold of 3 are re-routed to \texttt{Llama-3.3-70B} for a second opinion. We dropped this approach because the confidence scores turned out to be uncalibrated, in the wrong direction: accuracy conditional on a self-reported confidence of 3 was $0.713$, while accuracy conditional on the supposedly more certain confidence of 5 was only $0.500$, backwards from what a usable confidence signal would look like. Consistent with this, routing triggered on 244 of 250 paragraphs (97.6\%) because the small model rarely reported high confidence at all, and overall performance after routing (accuracy $0.708$, macro-$F_1 = 0.619$, $\kappa = 0.447$) was worse than \texttt{ministral-3-14b} annotating alone without any routing (accuracy $0.796$, macro-$F_1 = 0.618$, $\kappa = 0.601$; Table~\ref{tab:confidence}). We report this as a negative result: small LLMs cannot be assumed to reliably self-assess their own uncertainty, and that assumption should be tested, not taken for granted, before building a routing pipeline around it.

\subsubsection{Multi-Model Ensembling}
\label{sec:method-ensemble}

Since self-reported confidence failed as an uncertainty signal, we tried a more principled alternative: actual disagreement between two independently annotating models, which is an objective signal rather than a subjective self-assessment. We ran \texttt{ministral-3-14b} (Model A) and \texttt{Llama-3.3-70B}\footnote{Quantized variant \texttt{Meta-Llama-3.3-70B-Instruct-AWQ-INT4}, served via the university API.} (Model B) independently on the same test set; where the two models agreed, we kept their shared label; where they disagreed (74 of 250 paragraphs, 29.6\%), a third model, \texttt{qwen3.5-27b}, adjudicated the conflict by inspecting both labels and their explanations. This ensemble reached accuracy $0.800$ and macro-$F_1 = 0.684$ ($\kappa = 0.610$), and the tie-breaker sided with Model A on 61 of the 74 conflict paragraphs and Model B on 13, for a tie-break accuracy of $0.649$ against human ground truth (Table~\ref{tab:ensemble}). We nonetheless did not adopt this configuration, for two reasons. First, the improvement over \texttt{ministral-3-14b} alone was marginal ($\kappa = 0.610$ vs.\ $0.593$) relative to the cost of running two models plus a tiebreaker on every paragraph. Second, because this is the same test set used to develop the few-shot instructions in the revision round above, not an independently drawn one, even this marginal gain is subject to the overfitting risk discussed in Section~\ref{sec:limitations} and should not be read as a validated generalization estimate. Taken together, both comparisons pointed to the same conclusion: the best-performing configuration we found for this task was the single \texttt{ministral-3-14b} model with well-chosen few-shot examples (Section~\ref{sec:method}), not a multi-model pipeline, and that is the configuration we carried forward.

\subsection*{Step 4: Final Representative Validation}
\label{sec:method-final}

% TODO(authors): replace this subsection once the n=600 file is available. Suggested content:
% - how the 600-paragraph sample was drawn (representative vs. targeted; from which pool)
% - which instruction version / model configuration was used (final few-shot instructions? ensemble?)
% - final human intercoder Kappa/alpha = 0.76 and how it compares to the 250/new250 batches
% - final LLM validity metrics (accuracy, macro-F1, kappa/alpha vs. human ground truth) on this sample
Building on the improvements identified above, we carried out a final round of representative human validation on an independently drawn sample of 600 paragraphs, double-coded by both authors, and re-ran our best-performing LLM configuration on the same sample to obtain a final, single validity estimate that is not tied to either of the two smaller 250-paragraph test sets used for instruction development in Step 3 and the subsequent revision round. Reporting reliability and validity on a sample not used at any point for instruction tuning is important precisely because Section~\ref{sec:discussion} shows that validity measured on a repeatedly-tuned-against sample can overstate how well a pipeline generalizes; a fresh 600-paragraph sample is our attempt to give a less optimistic, more trustworthy final number. \todo{insert final sample composition, human intercoder Krippendorff's $\alpha = 0.76$, and LLM validity metrics (accuracy, macro-$F_1$, $\kappa$).}

\section{Results}
\label{sec:results}

\subsection{Human Intercoder Reliability}
\label{sec:results-reliability}

Table~\ref{tab:reliability} summarizes human intercoder reliability across the batches double-coded by both authors. Reliability rose substantially between the first 250-paragraph batch (used for Step 3 and the subsequent revision-round model validation) and the independently drawn new250 batch, reflecting cumulative codebook refinement: raw percentage agreement rose from 76.8\% to 94.4\%, and Krippendorff's $\alpha$ from 0.560 to 0.786. Using the interpretation bands common in this literature ($\alpha \geq 0.8$ very good, $0.6$--$0.8$ good, $0.4$--$0.6$ moderate, $0.2$--$0.4$ fair, below $0.2$ poor; \citealp{weberreichardt2023}), this represents a move from merely moderate agreement to good agreement, just short of the very-good threshold. This improvement in reliability coincided with a shift in label distribution: the first batch was comparatively balanced across labels (108 \texttt{NO\_CRIME\_FRAME}, 135 \texttt{CRIME\_FRAME\_NEG}, 7 \texttt{CRIME\_FRAME\_POS}), while the new250 batch skewed more heavily toward \texttt{NO\_CRIME\_FRAME} (213 of 250), which we discuss further in Section~\ref{sec:discussion}. Pooled across all 500 human-coded paragraphs from both batches, the final human label distribution was 69\% \texttt{NO\_CRIME\_FRAME}, 30\% \texttt{CRIME\_FRAME\_NEG}, and only 1\% \texttt{CRIME\_FRAME\_POS}, the empirical basis for collapsing the label scheme to a binary judgment (Section~\ref{sec:method}).

\begin{table}[htbp]
\centering
\begin{tabular}{lrrr}
\toprule
\textbf{Batch} & \textbf{n} & \textbf{\% Agreement} & \textbf{Krippendorff's $\alpha$} \\
\midrule
First batch (Step 3 / revision test set) & 250 & 76.8\% & 0.560 \\
new250 (independent batch)      & 250 & 94.4\% & 0.786 \\
Step 4: final validation sample & 600 & \todo{} & 0.76~\todo{confirm} \\
\bottomrule
\end{tabular}
\caption{Human intercoder reliability across annotation batches.}
\label{tab:reliability}
\end{table}

\subsection{LLM Validity Against Human Ground Truth}

Table~\ref{tab:llm-validity} reports LLM annotation validity, measured against the human consensus labels, across the successive stages and prompting conditions described in Section~\ref{sec:method}. Every row above the final validation sample was computed under the original three-class scheme (\texttt{NO\_CRIME\_FRAME}/\texttt{CRIME\_FRAME\_NEG}/\texttt{CRIME\_FRAME\_POS}), since these are the formative stages that led to the decision to collapse the scheme to binary (Section~\ref{sec:results-reliability}); only the final validation sample reflects the binary scheme actually used going forward. Within this formative, three-class evidence, the best result on any 250-paragraph test set was obtained with the revised few-shot instructions on the original revision-round test set (macro-$F_1 = 0.624$); the same instructions, applied to the independently drawn new250 set, achieved lower validity (macro-$F_1 = 0.492$) despite that set's higher human reliability.

\begin{table}[htbp]
\centering
\begin{tabular}{lrrrr}
\toprule
\textbf{Stage / condition} & \textbf{n} & \textbf{Accuracy} & \textbf{Macro-$F_1$} & \textbf{$\kappa$} \\
\midrule
Step 3 (curated hard/soft sample) & 250 & 0.656 & 0.521 & 0.321 \\
Revision: A. Zero-shot              & 250 & 0.644 & 0.492 & 0.304 \\
Revision: B. Few-shot                & 250 & 0.668 & 0.509 & 0.380 \\
Revision: C. Few-shot + CoT          & 250 & 0.660 & 0.502 & 0.351 \\
Revised instructions: B. Few-shot      & 250 & 0.792 & \textbf{0.624} & 0.588 \\
Revised instructions: C. Few-shot + CoT & 250 & 0.732 & 0.654 & 0.472 \\
new250 (vs.\ final human ground truth) & 250 & 0.804 & 0.492 & 0.498 \\
Step 4: final validation sample & 600 & \todo{} & \todo{} & \todo{} \\
\bottomrule
\end{tabular}
\caption{LLM annotation validity (\texttt{ministral-3-14b}) against human ground truth, by pipeline stage.}
\label{tab:llm-validity}
\end{table}

\subsection{Model Comparison: Ensembling vs.\ Confidence-Based Routing}

Table~\ref{tab:ensemble} and Table~\ref{tab:confidence} report the two side-experiments described in Sections~\ref{sec:method-ensemble} and~\ref{sec:method-confidence}. Multi-model ensembling with tie-break adjudication (macro-$F_1 = 0.684$) was the single best-performing configuration we tested anywhere in this project, improving substantially on any single-model result in Table~\ref{tab:llm-validity}. Confidence-based routing, in contrast, underperformed even simple single-model annotation, because the small model's self-reported confidence turned out not to track its actual accuracy.

\begin{table}[htbp]
\centering
\begin{tabular}{lrrr}
\toprule
\textbf{Configuration} & \textbf{Accuracy} & \textbf{Macro-$F_1$} & \textbf{$\kappa$} \\
\midrule
Ministral-3-14b alone & 0.792 & 0.615 & 0.593 \\
Ensemble (+ Llama-3.3-70B, Qwen3.5-27b tie-break) & 0.800 & \textbf{0.684} & 0.610 \\
\bottomrule
\end{tabular}
\caption{Two-model ensemble vs.\ single-model annotation ($n=250$; 74 tie-break cases, adjudicator sided with Model A on 61 and Model B on 13).}
\label{tab:ensemble}
\end{table}

\begin{table}[htbp]
\centering
\begin{tabular}{lrrr}
\toprule
\textbf{Configuration} & \textbf{Accuracy} & \textbf{Macro-$F_1$} & \textbf{$\kappa$} \\
\midrule
Ministral-3-14b alone (no routing) & 0.796 & 0.618 & 0.601 \\
Confidence-routed (244/250 rows routed) & 0.708 & 0.619 & 0.447 \\
\bottomrule
\end{tabular}
\caption{Confidence-based routing vs.\ single-model annotation ($n=250$).}
\label{tab:confidence}
\end{table}

\subsection{Substantive Framing Patterns Across Groups and Outlets}
\label{sec:results-substantive}

Answering our research question requires moving from pipeline validity to substantive measurement: applying the validated annotation configuration to a sample large and representative enough to compare groups and outlet categories against each other, rather than only against a single pooled accuracy figure. Following \citeauthor{freudenthaler2025}'s proposed use of the framework (their Figure~2), the natural way to summarize this is a per-group rate of \texttt{CRIME\_FRAME\_NEG} and \texttt{CRIME\_FRAME\_POS} (as a share of all paragraphs mentioning that group), compared against each group's overall share of corpus paragraphs, and cross-tabulated against the four outlet categories in Table~\ref{tab:outlets} to test whether the same group is framed differently in mainstream, regional, and alternative-right versus alternative-left coverage. A group whose \texttt{CRIME\_FRAME\_NEG} rate is high specifically in alternative-right outlets but not elsewhere, for instance, would indicate outlet-driven framing divergence rather than a property of the group's coverage in German news generally.

% TODO(authors): populate this subsection once the group x outlet-type crosstab has been computed from the n=600 (or larger) annotated sample.
% Suggested structure:
% - overall label distribution on the final annotated sample
% - which groups receive disproportionate CRIME_FRAME_NEG vs. CRIME_FRAME_POS relative to their share of paragraphs
% - whether this differs by outlet category (mainstream national / regional / alternative right / alternative left), per Table~\ref{tab:outlets}
% - a table or figure analogous to Freudenthaler et al.'s Figure 2 (mapping groups on a threatening-vs-beneficial axis)
\todo{report the group $\times$ outlet-category crosstab here, once the corresponding annotated sample is available: overall label distribution, which groups show disproportionate \texttt{CRIME\_FRAME\_NEG}/\texttt{CRIME\_FRAME\_POS} rates relative to their corpus share, and whether this varies across the four outlet categories in Table~\ref{tab:outlets}.}

\section{Discussion}
\label{sec:discussion}

The clearest pattern across our results is a dissociation between human intercoder reliability and LLM validity: reliability between our two human coders rose sharply between the first and second batches ($\alpha: 0.560 \rightarrow 0.786$) as we clarified genuinely ambiguous decision rules (most notably around suspicion/investigation language and the war-reporting exclusion), but LLM macro-$F_1$ against human ground truth did not rise correspondingly, and on the independently drawn new250 batch it fell back close to its original Step~3 level. We see at least two candidate explanations, which are not mutually exclusive. First, the revised instructions may have partly overfit to the specific error patterns of the first 250-paragraph test set (the few-shot examples were themselves written in response to errors observed there), so that a fresh sample surfaced new error patterns the instructions did not yet address: a risk inherent to iterating instructions against the same fixed sample repeatedly, and one we flag explicitly as a limitation (Section~\ref{sec:limitations}). Second, the new250 batch's more skewed label distribution (213/250 \texttt{NO\_CRIME\_FRAME}) is exactly the kind of class imbalance under which macro-averaged metrics are most sensitive to a model's performance on rare classes, and \texttt{CRIME\_FRAME\_POS} in particular remained a stubbornly low-support, low-recall class across every configuration we tested. This is not merely an artifact of our specific 250-paragraph test sets: \citet{shang2023} show formally that precision and recall estimated on a randomly sampled test set carry wide confidence intervals precisely when the positive class is rare, and that the width of those intervals shrinks only with deliberate oversampling of the minority class, which is the same logic behind our own oversampled hard-case and Step~3 test sets (Section~\ref{sec:data}). At $n=250$ with a class as rare as \texttt{CRIME\_FRAME\_POS} (roughly 1\% of paragraphs, Section~\ref{sec:results-reliability}), our recall estimates for that class should be read as suggestive rather than precise.

Substantively, ensembling two independently prompted LLMs with a third model as adjudicator was the most effective lever we found for improving validity on the test set available to us (macro-$F_1 = 0.684$, Table~\ref{tab:ensemble}), outperforming every single-model configuration, including the best-tuned few-shot instructions; because this test set is the same one used to develop those instructions, this result should be read as promising rather than confirmed (Section~\ref{sec:limitations}). This is broadly consistent with the intuition that two differently sized, differently trained models make partially independent errors, so that disagreement is itself an informative signal about which paragraphs are genuinely hard rather than about which model is simply worse. Confidence-based routing, by contrast, failed for a specific and diagnosable reason: \texttt{ministral-3-14b}'s self-reported confidence was almost never high, so the routing rule triggered on 97.6\% of paragraphs and provided almost no discriminative signal about where the model actually needed help. This is a useful negative result for future work using small LLMs as self-assessing annotators: self-reported confidence should be validated against actual accuracy before being trusted as a routing criterion, rather than assumed to be informative by construction.

It is also worth situating our prompting-based approach against the fine-tuned classifier approach used elsewhere in this seminar for a different sub-dimension of the same framework (stereotype content warmth/coldness, on the same underlying corpus). A fine-tuned transformer classifier requires assembling a dedicated, sufficiently large labeled training set before any evaluation is possible and, as that parallel project also found, performs poorly on rare classes under severe class imbalance (in that case, warm and cold labels each made up under 5\% of a 1{,}200-paragraph human-coded sample). Our few-shot prompting approach sidesteps the need for a large labeled training set entirely, at the cost of requiring careful, hand-written instructions and examples rather than learned parameters, and remains sensitive to the same rare-class problem in a different way: \texttt{CRIME\_FRAME\_POS} was consistently the hardest class for our LLM annotators to recall, much as the minority classes were for a fine-tuned classifier. Neither approach, in other words, escapes the fundamental difficulty that framing categories of interest to racism research are, empirically, rare events in general news coverage.

Finally, it is worth being precise about what a validated security-crime-frame label does and does not tell us about where a group's overall portrayal falls on \citeauthor{freudenthaler2025}'s spectrum from overtly racist to egalitarian discourse (Section~\ref{sec:theory}). A high rate of \texttt{CRIME\_FRAME\_NEG} for a given group is, on its own, evidence of one ingredient of a threatening portrayal, not proof that the group is being essentialized, homogenized, or othered in the stronger sense the framework distinguishes: a single frame family can co-occur with a merely stereotyping portrayal just as easily as with an overtly racist one, and disentangling the two requires the other five dimensions this report does not measure (Section~\ref{sec:theory}). We therefore read our results as one necessary input into that larger mapping, to be joined with the other dimensions measured elsewhere in this seminar, rather than as a stand-alone verdict on any single group's treatment in German news. % TODO(authors): once Section~\ref{sec:results-substantive} is filled in, add 1-2 paragraphs here connecting the substantive group/outlet findings back to the framing-effects literature in Section~\ref{sec:framing-lit} and to Freudenthaler et al.'s expected patterns of group portrayal.

\section{Limitations and Future Work}
\label{sec:limitations}

\textbf{Instruction iteration on a fixed test set.} As discussed in Section~\ref{sec:discussion}, our few-shot examples and decision-rule refinements were developed by inspecting errors on the same 250-paragraph Step 3 test set across multiple revision rounds, and both the ensembling and confidence-routing side experiments (Sections~\ref{sec:method-ensemble}--\ref{sec:method-confidence}) were also evaluated on this same set rather than on an independently drawn one. While this is consistent with the recommended routine \citep{freudenthalerbacchuss}, it risks overfitting to that specific sample's error patterns rather than to the underlying concept, and it means the ensembling result in particular has not been confirmed on data the pipeline was not already tuned against. Future work should re-evaluate both side experiments on a fresh, independently drawn sample before treating the ensemble configuration as the pipeline's final design, and future iterations more generally should draw a fresh sample between every round of instruction revision, rather than reusing the same 250 paragraphs throughout Step 3 and the subsequent revision rounds.

\textbf{Persistent rarity of \texttt{CRIME\_FRAME\_POS}.} Across every test set and prompting condition, positively valenced security-crime framing was both rare in the human-coded ground truth (1\% pooled across our 500 human-coded paragraphs, Section~\ref{sec:results-reliability}) and difficult for the LLM to recall reliably. Future work should consider deliberately oversampling candidate \texttt{CRIME\_FRAME\_POS} paragraphs (analogous to the zero-shot oversampling already used for hard-case testing, Section~\ref{sec:data}) specifically to build a larger, dedicated few-shot and evaluation set for this class; \citet{shang2023} provide a formal method for choosing an oversampling ratio that minimizes the confidence-interval width of precision and recall for a rare positive class, which would let future work state, rather than merely note, its uncertainty about \texttt{CRIME\_FRAME\_POS} performance specifically.

\textbf{Confidence-based routing was not diagnostic.} Our negative result on confidence-based routing (Section~\ref{sec:method-confidence}) should not be read as evidence against routing strategies in general, only against untested self-reported confidence as a routing signal. A calibration study, checking whether a model's self-reported confidence actually correlates with its accuracy before deploying it as a routing criterion, would be a natural precursor to any future routing pipeline.

\textbf{Scale of the final validated pipeline.} Our best-performing configuration (multi-model ensembling, Section~\ref{sec:method-ensemble}) was validated on 250 paragraphs and has not been applied to the full 659{,}895-paragraph corpus; a partial re-annotation attempt at scale (targeting paragraphs previously labeled as security-negative or -positive under an earlier instruction version) was left incomplete during this project. Extending ensemble annotation to the full corpus, or to a large representative sample drawn specifically for that purpose, is the clearest next step for turning the validated pipeline into a corpus-wide substantive analysis of the kind sketched in Section~\ref{sec:results-substantive}.

\textbf{Single frame, single country, single year.} We measure one frame family (security-crime) within one of \citeauthor{freudenthaler2025}'s six proposed dimensions, in German news from a single year (2022). The framework itself is explicitly designed to support comparison across dimensions, time periods, and national contexts \citep{freudenthaler2025}; extending this pipeline to additional years, additional frame families (economic and cultural threat/benefit), or additional national media systems would let future work test whether the patterns we observe here are stable properties of German news discourse or specific to this particular year and frame.

\section{Conclusion}
\label{sec:conclusion}

This project developed and transparently documented a four-stage LLM annotation pipeline for measuring the security-crime frame applied to ethnic, national, and religious groups in German news coverage, grounded in \citeauthor{freudenthaler2025}'s framework for measuring fragmented racism in large-scale corpora. We found that human intercoder reliability and LLM validity did not necessarily move together across iterations, that multi-model ensembling meaningfully outperformed both single-model annotation and confidence-based routing, and that confidence-based routing failed specifically because a small model's self-reported confidence was not a reliable signal of its own accuracy. % TODO: add one sentence on the headline substantive finding once Section 5.4 is complete. Beyond the specific numbers reported here, we hope the level of methodological detail in Sections~\ref{sec:method} and~\ref{sec:discussion}, including the stages that did not work as expected, is itself a useful contribution for other groups building LLM-assisted content-analysis pipelines on top of the same framework.

\subsection*{Data and Code Availability}
The raw news corpus was provided by the course instructor for use within this seminar and is not publicly redistributable. The annotation instructions (Appendix~\ref{app:instructions}), analysis code, and derived, non-identifying result files (label counts, agreement statistics, evaluation summaries) referenced throughout this report are available in the project repository on request from the authors.

\newpage
\bibliography{references}

\newpage
\appendix

\section{Appendix}

\subsection{Annotation Instructions (Full Text)}
\label{app:instructions}

The instructions below are the final version used throughout Steps 2--4 (Section~\ref{sec:method}), reproduced verbatim (German original) from the annotation pipeline code, together with the reminder text appended to every prompt.

\begin{instructionbox}
\textbf{System role:} \enquote{Sie sind ein Experte f\"ur Inhaltsanalyse. Befolgen Sie die Annotationsanweisungen genau und antworten Sie immer im angegebenen Format.}

\medskip
\textbf{Instructions (excerpt, decision rules):}

\texttt{NO\_CRIME\_FRAME} is assigned when no explicit crime or security frame is present, in particular when the group is not mentioned; crime is mentioned but not linked to the group; the group is portrayed as a victim rather than a perpetrator; crime is attributed to social conditions rather than the group; the text relativizes or contradicts a crime attribution (\enquote{kein Generalverdacht}, \enquote{nicht alle}); war, geopolitical conflict, or military events are described; humanitarian crises or natural disasters are described; a perpetrator is mentioned who is \emph{not} a member of the group; fictional, cultural, or satirical depictions are involved; or demonstrations/protests are mentioned without a link to crime, violence, or concrete security measures.

\texttt{CRIME\_FRAME\_NEG} is assigned when the group is portrayed as a source of crime or security threat: as perpetrator, suspect, or accused; linked to concrete offenses; portrayed as a cause of rising crime or a danger to public or national security; linked to terrorism or extremism; or portrayed as endorsing, legitimizing, celebrating, or supporting an attack or act of violence \emph{even without direct involvement}. Suspicion, investigation, and as-yet-unconfirmed allegations (\enquote{mutmaßlich}, \enquote{verdächtigt}, \enquote{Ermittlungen gegen}) count as \texttt{CRIME\_FRAME\_NEG} once explicitly linked to the group; a judicial conviction is not required. \textbf{Exception:} the label \enquote{illegale Einwanderer} (illegal immigrants) does \emph{not}, by itself, count toward this label unless a further offense is named; smugglers (\enquote{Schleuser}) from the group described as criminals do count.

\texttt{CRIME\_FRAME\_POS} is assigned when a security-crime frame is present but portrayed positively: crime-prevention measures, prevention work (social work, integration, deradicalization), explicit differentiation from generalized suspicion, crime explained via social factors rather than the group itself, or group members actively contributing to safety or prevention (e.g.\ police, social work). Demonstrations count as \texttt{CRIME\_FRAME\_POS} only when explicitly tied to concrete legal, political, or preventive security changes.

\medskip
\textbf{Fallback rule (\enquote{Im Zweifelsfall}):} \enquote{Ask yourself whether a reader would perceive the group as criminal, suspect, security-threatening, supportive of violence, or as a safety factor. If not $\rightarrow$ \texttt{NO\_CRIME\_FRAME}.}

\medskip
\textbf{Reminder (appended to every prompt):} \enquote{Vergeben Sie genau ein Label: NO\_CRIME\_FRAME, CRIME\_FRAME\_NEG oder CRIME\_FRAME\_POS. \dots\ Kriegskontexte, fiktionale Inhalte und blo{\ss}e Demonstrationen ohne Sicherheits- oder Pr\"aventionsbezug sind NO\_CRIME\_FRAME.}
\end{instructionbox}

\subsection{Label Scheme Summary}
\label{app:codebook}

\begin{table}[htbp]
\centering
\begin{tabularx}{\textwidth}{lX}
\toprule
\textbf{Label} & \textbf{Definition} \\
\midrule
\texttt{NO\_CRIME\_FRAME} & No explicit link between the group and crime, security, danger, violence, terrorism, suspicion, investigation, or prevention; includes victim portrayals, war/geopolitical reporting, and fictional depictions. \\
\texttt{CRIME\_FRAME\_NEG} & Group explicitly portrayed as source of crime, danger, or security threat (perpetrator, suspect, accused, endorser of violence). \\
\texttt{CRIME\_FRAME\_POS} & Security-crime frame present but positively valenced: group contributes to safety, prevention, or is explicitly differentiated from generalized suspicion. \\
\bottomrule
\end{tabularx}
\caption{Summary of the final annotation label scheme (Appendix~\ref{app:instructions} gives the full instructions).}
\end{table}

\clearpage
\section*{Ehrenw\"ortliche Erkl\"arung}

Ich versichere, dass ich die beiliegende Bachelor-, Master-, Seminar-, oder Projektarbeit ohne Hilfe Dritter und ohne Benutzung anderer als der angegebenen Quellen und in der untenstehenden Tabelle angegebenen Hilfsmittel angefertigt und die den benutzten Quellen w\"ortlich oder inhaltlich entnommenen Stellen als solche kenntlich gemacht habe. Diese Arbeit hat in gleicher oder \"ahnlicher Form noch keiner Pr\"ufungsbeh\"orde vorgelegen. Ich bin mir bewusst, dass eine falsche Erkl\"arung rechtliche Folgen haben wird.

% TODO(authors): fill in this table yourselves; see chat for context on why this was left as a placeholder.
\begin{center}
  \textbf{Declaration of Used AI Tools} \\[.3em]
  \begin{tabularx}{\textwidth}{lXlc}
    \toprule
    Tool & Purpose & Where? & Useful? \\
    \midrule
    \todo{} & \todo{} & \todo{} & \todo{} \\
    \bottomrule
  \end{tabularx}
\end{center}

\vspace{2cm}
\noindent Robin-Kiara Braun, Marmee Pandya \\
\noindent Mannheim, den {\today}

\end{document}
